\begin{summarybox}

	\textbf{\textcolor{CSummary}{一句话核心}}

	\textbf{当三角形内一条平行于底边的直线截两边，所得小三角形与原三角形相似，并由此得到一系列比例关系。}

	\medskip
	\textbf{\textcolor{CSummary}{使用条件}}
	\begin{itemize}[leftmargin=*, itemsep=0.5em, topsep=0.4em]
		\item \textbf{条件1（平行线条件）：}$DE\parallel BC$，且 $D$ 在 $AB$ 上，$E$ 在 $AC$ 上。
		\item \textbf{条件2（对应点位置）：} 小三角形顶点 $D,E$ 分别在大三角形边 $AB,AC$ 上。
	\end{itemize}

	\medskip
	\textbf{\textcolor{CSummary}{关键提醒}}
	\begin{itemize}[leftmargin=*, itemsep=0.5em, topsep=0.4em]
		\item \textbf{易错点（面积比平方）：} 面积比是相似比的平方，不能直接用边长比代替面积比。
		\item \textbf{检查项（对应边顺序）：} 写比例时必须严格按顶点对应，避免 $\frac{AD}{DB}=\frac{DE}{BC}$ 的错误。
	\end{itemize}
\end{summarybox}
